\documentclass[11pt]{article}
\usepackage[T1]{fontenc}
\usepackage[utf8]{inputenc}
\usepackage{lmodern}
\usepackage[margin=1in]{geometry}
\usepackage{amsmath,amssymb,amsthm,mathtools}
\usepackage{enumitem}
\usepackage[hidelinks]{hyperref}

\newtheorem{theorem}{Theorem}[section]
\newtheorem{proposition}[theorem]{Proposition}
\newtheorem{definition}[theorem]{Definition}
\newtheorem{lemma}[theorem]{Lemma}
\newtheorem{corollary}[theorem]{Corollary}
\newtheorem{remark}[theorem]{Remark}

\title{Complete Commutator-Word Tomography of Alternating Pencils}
\author{Luca Blanchi}
\date{June 2026}

\begin{document}
\maketitle

\begin{abstract}
Let $k=\mathbb F_p$, $p\neq2$, and let

\[
\beta:\Lambda^2_k V\to k^2
\]

be an alternating pencil. After choosing coordinates on $k^2$, write

\[
\beta=(A,B),
\]

where $A,B\in\Lambda^2V^*$. Surjective radical-free alternating pencils are precisely the commutator tensors of special class-two exponent-$p$ groups whose derived subgroup has dimension $2$. More generally, every alternating pencil defines a class-two exponent-$p$ group, and the special condition corresponds to surjectivity and zero common radical.

We prove that, for fixed $k$ and $N=\dim V$, there is an explicit finite family of Fourier-tagged commutator-word distributions whose values determine the full skew-symmetric Kronecker canonical data of the pencil. The singular part is recovered from rectangular chain probes. If $m_e$ denotes the number of singular Kronecker blocks of minimal index $e$, then the chain defects

\[
\delta_f
=
f\dim V-\operatorname{rank}(F_f\otimes A+G_f\otimes B)
\]

satisfy

\[
\delta_f=\sum_{e\ge0}m_e\max(0,f-e),
\]

and hence

\[
m_e=\delta_{e+1}-2\delta_e+\delta_{e-1}.
\]

The regular part is recovered from closed-point Artin probes. For a regular hyperbolic block with elementary divisor ($\xi,e$), and a closed-point Artin probe ($\eta,f$), we prove the local block-pairing formula

\[
\Delta((\xi,e),(\eta,f))
=
2\deg(\xi)\min(e,f)\mathbf 1_{\xi=\eta}.
\]

We also compute explicitly the contribution of singular Kronecker blocks to these Artin probes:

\[
\operatorname{rad}_{\eta,f}^{\mathrm{sing}}
=
M_{\mathrm{sing}}\deg(\eta)f,
\]

where $M_{\mathrm{sing}}$ is the total number of singular Kronecker blocks. Since the singular part is already recovered by chain probes, this contribution can be subtracted, leaving the regular divisor multiplicities.

Combining these ingredients, the Fourier-tagged commutator-word distributions recover the singular minimal-index multiset and the regular elementary-divisor divisor, up to the natural $\operatorname{PGL}_2(k)$-action on the parameter line. Consequently, for special class-two exponent-$p$ groups $G_\beta$ with

\[
\dim_{\mathbb F_p}G'=2,
\]

these distributions determine the group isomorphism class.

The result is exact and finite. It is not presented as a practical isomorphism algorithm; the contribution is the complete word-moment observability of the Kronecker canonical data of an alternating pencil.

\end{abstract}

\section*{Declaration of generative AI and AI-assisted technologies in the writing process}
This article belongs to the Presentation theory section. Generative AI, including ChatGPT by \mbox{OpenAI}, was used to assist with mathematical drafting, formalization, review, and editing. The note may be preliminary, may have been only partially reviewed, and in some cases may be only AI-reviewed.

\section{Introduction}

Let $p$ be an odd prime and let $k=\mathbb F_p$. A special finite $p$-group $G$ of nilpotency class two and exponent $p$ has

\[
V=G/Z(G),
\qquad
W=G'=Z(G),
\]

and its commutator induces an alternating bilinear map

\[
\beta:\Lambda^2_k V\to W.
\]

Conversely, every such alternating map gives a class-two exponent-$p$ group

\[
G_\beta=V\oplus W
\]

with multiplication

\[
(v,z)(v',z')
\left(v+v',
z+z'+\frac12\beta(v,v')\right).
\]

The commutator is

\[
[(v,z),(v',z')]=(0,\beta(v,v')).
\]

When $\dim W=2$, the commutator tensor is an alternating pencil. After choosing a basis of $W$, write

\[
\beta=(A,B),
\]

where $A,B$ are alternating forms on $V$. The corresponding matrix pencil is

\[
A+tB.
\]

Alternating pencils over fields of characteristic different from $2$ have a classical canonical decomposition into regular elementary-divisor blocks and singular Kronecker blocks. This is the skew-symmetric version of Kronecker pencil theory, and it appears in canonical-form treatments of pairs of forms and pairs of skew-symmetric matrices under congruence.

The purpose of this paper is to show that the entire canonical datum of an alternating pencil is observable from commutator-word distributions.

The guiding mechanism is Fourier-tagged word tomography. Given rectangular coefficient matrices $C,D$, one forms central commutator outputs whose Fourier coefficients are

\[
p^{-\operatorname{rank}(C\otimes A+D\otimes B)}.
\]

Thus commutator-word distributions can be used as rank probes for the pencil.

The singular Kronecker blocks are detected by rectangular chain probes. The regular elementary divisors are detected by closed-point Artin probes, which evaluate the pencil at infinitesimal neighbourhoods of closed points of $\mathbb P^1$. Together these probes recover the full canonical datum.

\subsection{Main theorem}

The main theorem is the following.

\begin{theorem}[Theorem A -- Complete commutator-word tomography of alternating pencils]

Fix a finite field $k=\mathbb F_p$, $p\neq2$, and an integer $N$. There exists an explicit finite family of Fourier-tagged commutator-word distributions such that, for every alternating pencil

\[
\beta:\Lambda^2_k V\to k^2,
\qquad
\dim V=N,
\]

these distributions determine the projective congruence class of $\beta$.

Equivalently, they determine:

\begin{itemize}
\item the multiplicities of all singular Kronecker blocks;
\end{itemize}

\begin{itemize}
\item the multiplicities of all regular elementary divisors ($\xi,e$), up to the natural $\operatorname{PGL}_2(k)$-action on $\mathbb P^1$.
\end{itemize}

Consequently, for special class-two exponent-$p$ groups $G_\beta$ with

\[
\dim_{\mathbb F_p}G'=2,
\]

the same finite family determines the group isomorphism class.

\end{theorem}

\subsection{Scope}

The theorem is exact and finite. It is not a sampling theorem and not an optimized isomorphism algorithm. The orbit-coding step used to pass from a chosen target basis to an intrinsic projective invariant is finite but not intended to be computationally efficient.

The contribution is instead structural: the complete skew-symmetric Kronecker canonical data of an alternating commutator pencil is observable from explicitly designed commutator-word distributions.

\section{Class-two exponent-$p$ groups and alternating pencils}

Let $k=\mathbb F_p$, $p\neq2$, and let

\[
\beta:\Lambda^2_k V\to W
\]

be alternating.

Define

\[
G_\beta=V\oplus W
\]

with multiplication

\[
(v,z)(v',z')
\left(
v+v',
z+z'+\frac12\beta(v,v')
\right).
\]

Then $G_\beta$ is a group of exponent $p$ and nilpotency class at most two. Its commutator is

\[
[(v,z),(v',z')]
(0,\beta(v,v')).
\]

If $\beta$ is surjective and has zero radical,

\[
\operatorname{rad}\beta
=
\{v\in V:\beta(v,V)=0\}=0,
\]

then

\[
G_\beta'=Z(G_\beta)=W,
\]

so $G_\beta$ is special.

Conversely, every special class-two exponent-$p$ group arises this way.

Two tensors

\[
\beta:\Lambda^2V\to W,
\qquad
\beta':\Lambda^2V'\to W'
\]

are pseudo-isometric if there are vector-space isomorphisms

\[
S:V\to V',
\qquad
T:W\to W'
\]

such that

\[
T\beta(v,v')=\beta'(Sv,Sv')
\]

for all $v,v'\in V$. For special class-two exponent-$p$ groups, group isomorphism is equivalent to pseudo-isometry of the commutator tensor.

Throughout the main body of this paper,

\[
W=k^2.
\]

After choosing a basis of $W$, we write

\[
\beta=(A,B),
\]

where

\[
A,B\in\Lambda^2V^*.
\]

Changing the basis of $W$ acts on the projective parameter line by $\operatorname{PGL}_2(k)$.

\section{Fourier-tagged commutator distributions}

Let

\[
C,D\in M_{r\times s}(k).
\]

For variables

\[
x_1,\ldots,x_s,\qquad y_1,\ldots,y_r
\]

in $G_\beta$, define two central outputs

\[
Z_C
=
\sum_{i=1}^{r}\sum_{j=1}^{s} C_{ij}\,\beta(\bar x_j,\bar y_i),
\]

\[
Z_D
=
\sum_{i=1}^{r}\sum_{j=1}^{s} D_{ij}\,\beta(\bar x_j,\bar y_i),
\]

where $\bar x_j,\bar y_i$ denote the images in $V=G/Z(G)$.

Thus

\[
(Z_C,Z_D)
\]

is a random variable in

\[
W\times W.
\]

Let

\[
P_{C,D}
\]

be its probability distribution.

Let

\[
\psi:k\to\mathbb C^\times
\]

be a nontrivial additive character. For

\[
\alpha,\gamma\in W^*,
\]

define the Fourier coefficient

\[
\widehat P_{C,D}(\alpha,\gamma)
=
\mathbb E\,
\psi\bigl(\alpha(Z_C)+\gamma(Z_D)\bigr).
\]

For $\ell\in W^*$, write

\[
A_\ell=\ell\circ\beta\in\Lambda^2V^*.
\]

If $W=k e_A\oplus k e_B$ and

\[
e_A^\circ\circ\beta=A,\qquad e_B^\circ\circ\beta=B,
\]

then

\[
A_{e_A^\circ}=A,\qquad A_{e_B^\circ}=B.
\]

\begin{lemma}[Bilinear character average]

Let

\[
B_0:U_1\times U_2\to k
\]

be a bilinear form between finite-dimensional $k$-vector spaces. Then

\[
\mathbb E_{u\in U_1,\ v\in U_2}
\psi(B_0(u,v))
=
p^{-\operatorname{rank}B_0}.
\]

\end{lemma}

\begin{proof}

For fixed $u$, the inner average over $v$ equals $1$ if the linear functional

\[
v\mapsto B_0(u,v)
\]

is zero, and equals $0$ otherwise. Thus the average is the proportion of $u\in U_1$ lying in the left radical of $B_0$. That radical has codimension $\operatorname{rank}B_0$, giving the formula.

\end{proof}

\begin{theorem}[Rectangular Fourier-tagged moment formula]

For all

\[
C,D\in M_{r\times s}(k)
\]

and all

\[
\alpha,\gamma\in W^*,
\]

one has

\[
\boxed{
\widehat P_{C,D}(\alpha,\gamma)
=
p^{-\operatorname{rank}(C\otimes A_\alpha+D\otimes A_\gamma)}.
}
\]

Here

\[
C\otimes A_\alpha+D\otimes A_\gamma
\]

is regarded as a bilinear form on

\[
(k^s\otimes V)\times(k^r\otimes V).
\]

In particular, in the chosen coordinates on $W=k^2$,

\[
\boxed{
\widehat P_{C,D}(e_A^\circ,e_B^\circ)
=
p^{-\operatorname{rank}(C\otimes A+D\otimes B)}.
}
\]

\end{theorem}

\begin{proof}

We compute

\[
\alpha(Z_C)+\gamma(Z_D)
=
\sum_{i,j}
C_{ij}A_\alpha(\bar x_j,\bar y_i)
+
\sum_{i,j}
D_{ij}A_\gamma(\bar x_j,\bar y_i).
\]

This is precisely the bilinear form

\[
C\otimes A_\alpha+D\otimes A_\gamma
\]

on

\[
(k^s\otimes V)\times(k^r\otimes V).
\]

The variables ($\bar x_j$) and ($\bar y_i$) are uniformly distributed in $V^s$ and $V^r$. Applying Lemma 3.1 gives

\[
\widehat P_{C,D}(\alpha,\gamma)
=
p^{-\operatorname{rank}(C\otimes A_\alpha+D\otimes A_\gamma)}.
\]

\end{proof}

This is the basic interface between commutator-word distributions and the linear algebra of matrix pencils.

\section{Canonical data of alternating pencils}

We recall the canonical data used below.

A skew-symmetric pencil

\[
A+tB
\]

over a field of characteristic not (2) decomposes as a direct sum of two types of indecomposable blocks.

\subsection{Regular blocks}

The regular part is described by elementary divisors

\[
(\xi,e),
\]

where $\xi$ is a closed point of $\mathbb P^1_k$ and $e\ge1$. We encode the regular part by a decorated divisor

\[
\mathcal D_{\mathrm{reg}}
=
\sum_{\xi,e}m_{\xi,e}[\xi,e].
\]

A block $(\xi,e)$ contributes dimension

\[
2e\deg(\xi)
\]

to $V$.

\subsection{Singular Kronecker blocks}

The singular part is described by minimal indices

\[
e\ge0.
\]

Let $m_e$ be the multiplicity of the singular Kronecker block of index $e$. Such a block has dimension

\[
2e+1.
\]

The full canonical data is

\[
\left(
\mathcal D_{\mathrm{reg}},
(m_e)_{e\ge0}
\right),
\]

up to the natural $\operatorname{PGL}_2(k)$-action on the parameter line $\mathbb P^1$.

The existence and uniqueness of this canonical data are classical; in this paper the canonical form is used as an input from the theory of pairs of skew-symmetric forms under congruence.

\section{Singular chain probes}

We first recover the singular Kronecker minimal indices.

For $f\ge1$, define

\[
F_f,G_f\in M_{f\times(f+1)}(k)
\]

by

\[
(F_f)_{i,i}=1,
\qquad
(G_f)_{i,i+1}=1,
\qquad
0\le i\le f-1,
\]

and all other entries zero.

Given a pencil $(A,B)$, define

\[
\mathcal L_f(A,B):V^{f+1}\to(V^*)^f
\]

by

\[
\mathcal L_f(A,B)(z_0,\ldots,z_f)
=
(Az_0+Bz_1,
Az_1+Bz_2,
\ldots,
Az_{f-1}+Bz_f).
\]

Then

\[
\operatorname{rank}\mathcal L_f(A,B)
=
\operatorname{rank}(F_f\otimes A+G_f\otimes B).
\]

Define the chain defect

\[
\boxed{
\delta_f(A,B)
=
f\dim V-\operatorname{rank}(F_f\otimes A+G_f\otimes B).
}
\]

Equivalently,

\[
\delta_f(A,B)
=
\dim\ker \mathcal L_f(A,B)-\dim V,
\]

because

\[
\dim V^{f+1}=(f+1)\dim V.
\]

By Theorem 3.2, the value

\[
\operatorname{rank}(F_f\otimes A+G_f\otimes B)
\]

is recovered from a Fourier coefficient of an explicit commutator-word distribution.

\begin{lemma}[Regular blocks have zero chain defect]

If $(A,B)$ is a regular block, then

\[
\delta_f(A,B)=0
\]

for every $f\ge1$.

\end{lemma}

\begin{proof}

For a regular block, the pencil $A+tB$ is nonsingular over $k(t)$. After extending scalars and applying a projective change of parameter, we may assume that $A$ is invertible.

Then the equations

\[
Az_i+Bz_{i+1}=w_i,\qquad 0\le i\le f-1,
\]

can be solved recursively after choosing $z_f$. Hence

\[
\mathcal L_f(A,B):V^{f+1}\to(V^*)^f
\]

is surjective. Thus

\[
\operatorname{rank}\mathcal L_f=f\dim V,
\]

so

\[
\delta_f=0.
\]

\end{proof}

\subsection{Singular Kronecker blocks}

For $e\ge0$, define the rectangular pencil

\[
L_e(t):k^{e+1}\to k^e
\]

by

\[
L_e(t)(x_0,\ldots,x_e)
=
(x_0+t x_1,
x_1+t x_2,
\ldots,
x_{e-1}+t x_e).
\]

Write

\[
L_e(t)=F_e+tG_e.
\]

The corresponding skew-symmetric singular block is

\[
\mathcal K_e(t)
=
\begin{pmatrix}
0&L_e(t)^{\mathsf T}\\
-L_e(t)&0
\end{pmatrix}
\]

on

\[
V_e=k^{e+1}\oplus k^e.
\]

Thus

\[
\dim V_e=2e+1.
\]

\begin{lemma}[Singular chain defect]

For the singular block $\mathcal K_e$,

\[
\boxed{
\delta_f(\mathcal K_e)=\max(0,f-e).
}
\]

\end{lemma}

\begin{proof}

Write

\[
V_e=X\oplus Y,
\qquad
X=k^{e+1},
\qquad
Y=k^e.
\]

The chain equations for the $X$-part are

\[
F_ex_i+G_ex_{i+1}=0,
\qquad
0\le i\le f-1.
\]

In coordinates

\[
x_i=(x_{i,0},\ldots,x_{i,e}),
\]

these equations are

\[
x_{i,j}+x_{i+1,j+1}=0,
\qquad
0\le i\le f-1,\quad 0\le j\le e-1.
\]

The variables are identified along diagonals (j-i). There are

\[
f+e+1
\]

such diagonals. Hence the $X$-solution space has dimension

\[
f+e+1.
\]

The chain equations for the $Y$-part are dual:

\[
F_e^Ty_i+G_e^Ty_{i+1}=0,
\qquad
0\le i\le f-1.
\]

Writing

\[
y_i=(y_{i,0},\ldots,y_{i,e-1}),
\]

these equations are

\[
y_{i,0}=0,
\]

\[
y_{i,j}+y_{i+1,j-1}=0
\quad
(1\le j\le e-1),
\]

\[
y_{i+1,e-1}=0.
\]

They force all entries except a boundary family of dimension

\[
\max(0,e-f).
\]

Thus the $Y$-solution space has dimension

\[
\max(0,e-f).
\]

Therefore

\[
\dim\ker\mathcal L_f(\mathcal K_e)
=
(f+e+1)+\max(0,e-f).
\]

If $f<e$, this equals

\[
2e+1=\dim V_e.
\]

If $f\ge e$, this equals

\[
f+e+1
=
(2e+1)+(f-e).
\]

Hence

\[
\dim\ker\mathcal L_f(\mathcal K_e)-\dim V_e
=
\max(0,f-e).
\]

This is exactly

\[
\delta_f(\mathcal K_e).
\]

\end{proof}

\begin{theorem}[Recovery of singular minimal indices]

Let $m_e$ be the number of singular Kronecker blocks of minimal index $e$. Then

\[
\boxed{
\delta_f
=
\sum_{e\ge0}m_e\max(0,f-e).
}
\]

Consequently,

\[
\boxed{
m_e=\delta_{e+1}-2\delta_e+\delta_{e-1},
}
\]

with the convention

\[
\delta_0=\delta_{-1}=0.
\]

\end{theorem}

\begin{proof}

The chain defect is additive under direct sums. Regular blocks contribute zero by Lemma 5.1. Singular blocks contribute $\max(0,f-e)$ by Lemma 5.2. Hence

\[
\delta_f
=
\sum_{e\ge0}m_e\max(0,f-e).
\]

Taking first differences gives

\[
\delta_f-\delta_{f-1}
=
\sum_{e<f}m_e.
\]

Indeed, the summand $\max(0,f-e)$ increases by $1$ exactly when $e<f$. Taking a second difference yields

\[
(\delta_{e+1}-\delta_e)-(\delta_e-\delta_{e-1})
=
m_e.
\]

Thus

\[
\boxed{
m_e=\delta_{e+1}-2\delta_e+\delta_{e-1}.
}
\]

\end{proof}

Thus the singular Kronecker data is recovered from the Fourier-tagged rectangular distributions associated to the matrices $(F_f,G_f)$.

\section{Closed-point Artin probes for the regular part}

We now recover the regular elementary divisors.

Let $\xi$ be a closed point of $\mathbb A^1_k$, represented by an irreducible monic polynomial

\[
\varphi_\xi(t)\in k[t].
\]

Let

\[
d_\xi=\deg(\xi),
\qquad
K_\xi=k[t]/(\varphi_\xi),
\qquad
\alpha_\xi=t\bmod\varphi_\xi.
\]

For $e\ge1$, define

\[
M_{\xi,e}=k[t]/(\varphi_\xi^e),
\]

and let

\[
T_{\xi,e}:M_{\xi,e}\to M_{\xi,e}
\]

be multiplication by $t$.

The regular hyperbolic block associated to $(\xi,e)$ is

\[
H_{\xi,e}=M_{\xi,e}\oplus M_{\xi,e}^\vee.
\]

It carries the alternating forms

\[
A_{\xi,e}((x,\alpha),(y,\beta))
=
\alpha(y)-\beta(x),
\]

\[
B_{\xi,e}((x,\alpha),(y,\beta))
\alpha(T_{\xi,e}y)-\beta(T_{\xi,e}x).
\]

Equivalently,

\[
B_{\xi,e}
A_{\xi,e}(S_{\xi,e}\cdot,\cdot),
\]

where

\[
S_{\xi,e}=T_{\xi,e}\oplus T_{\xi,e}^\vee.
\]

Let $\eta$ be another closed point and define

\[
R_{\eta,f}=K_\eta[\pi]/(\pi^f),
\qquad
\theta_{\eta,f}=\alpha_\eta+\pi\in R_{\eta,f}.
\]

The closed-point Artin probe evaluates the pencil at $\theta_{\eta,f}$. On

\[
H_{\xi,e}\otimes_k R_{\eta,f}
\]

it gives the alternating form

\[
\Omega_{\xi,e;\eta,f}
B_{\xi,e}\otimes1-\theta_{\eta,f}(A_{\xi,e}\otimes1).
\]

Since $A_{\xi,e}\otimes1$ is nondegenerate,

\[
\operatorname{rad}\Omega_{\xi,e;\eta,f}
=
\ker(S_{\xi,e}\otimes1-\theta_{\eta,f}I).
\]

\begin{theorem}[Regular closed-point block pairing]

For regular hyperbolic blocks and closed-point Artin probes,

\[
\boxed{
\dim_k\operatorname{rad}\Omega_{\xi,e;\eta,f}
2\deg(\xi)\min(e,f)\mathbf 1_{\xi=\eta}.
}
\]

\end{theorem}

\begin{proof}

Because

\[
S_{\xi,e}=T_{\xi,e}\oplus T_{\xi,e}^\vee,
\]

the kernel of

\[
S_{\xi,e}\otimes1-\theta I
\]

is the direct sum of the corresponding kernel for $T_{\xi,e}$ and for its dual. These two kernels have the same $k$-dimension. Thus it is enough to compute

\[
\dim_k\ker(t-\alpha_\eta-\pi)
\]

on

\[
M_{\xi,e}\otimes_kR_{\eta,f}.
\]

If $\xi\neq\eta$, then $\varphi_\xi$ and $\varphi_\eta$ are coprime. Modulo $\pi$, the element

\[
t-\alpha_\eta
\]

is invertible in

\[
K_\eta[t]/(\varphi_\xi(t)^e).
\]

Therefore

\[
t-\alpha_\eta-\pi
\]

is invertible over the local Artin ring $R_{\eta,f}$. The kernel is zero.

Now assume $\xi=\eta$. Write

\[
K=K_\xi,
\qquad
\alpha=\alpha_\xi,
\qquad
R=K[\pi]/(\pi^f).
\]

After tensoring with $K$, only the local factor corresponding to the root $\alpha$ contributes to the kernel; the conjugate factors are invertible. On the contributing local factor, set

\[
u=t-\alpha.
\]

The local algebra is

\[
K[u,\pi]/(u^e,\pi^f),
\]

and the operator is multiplication by

\[
u-\pi.
\]

Since this is a linear operator on a finite-dimensional $K$-vector space, kernel and cokernel have the same $K$-dimension. The cokernel is

\[
K[u,\pi]/(u^e,\pi^f,u-\pi)
\cong
K[\pi]/(\pi^{\min(e,f)}).
\]

Thus the $K$-dimension of the kernel is

\[
\min(e,f).
\]

Multiplying by $[K:k]=\deg(\xi)$ gives

\[
\deg(\xi)\min(e,f).
\]

Doubling for the hyperbolic dual part gives

\[
2\deg(\xi)\min(e,f).
\]

\end{proof}

The point at infinity is handled by replacing $t$ by a local affine coordinate at infinity. The formula is invariant under $\operatorname{PGL}_2(k)$.

\section{Singular contribution to Artin probes}

To recover the regular part, we must subtract the known contribution of the singular blocks to the Artin probes. We record this explicitly.

Let

\[
L_e(t):k^{e+1}\to k^e
\]

be the rectangular Kronecker pencil

\[
L_e(t)(x_0,\ldots,x_e)
=
(x_0+t x_1,\ldots,x_{e-1}+t x_e),
\]

and let

\[
\mathcal K_e(t)
=
\begin{pmatrix}
0&L_e(t)^{\mathsf T}\\
-L_e(t)&0
\end{pmatrix}
\]

be the corresponding skew-symmetric singular block.

Let $\eta$ be a closed point and let

\[
R_{\eta,f}=K_\eta[\pi]/(\pi^f),
\qquad
\theta_{\eta,f}=\alpha_\eta+\pi.
\]

Evaluate $\mathcal K_e(t)$ at

\[
t=\theta_{\eta,f}.
\]

\begin{lemma}[Singular Artin contribution]

For every singular Kronecker block $\mathcal K_e$, the radical of the evaluated Artin form

\[
\mathcal K_e(\theta_{\eta,f})
\]

has $k$-dimension

\[
\boxed{
\deg(\eta)f.
}
\]

In particular, this contribution is independent of the minimal index $e$.

\end{lemma}

\begin{proof}

Write

\[
R=R_{\eta,f},
\qquad
\theta=\theta_{\eta,f}.
\]

The evaluated rectangular map is

\[
L_e(\theta):R^{e+1}\to R^e,
\]

\[
L_e(\theta)(x_0,\ldots,x_e)
=
(x_0+\theta x_1,\ldots,x_{e-1}+\theta x_e).
\]

This map is surjective. Given

\[
(y_0,\ldots,y_{e-1})\in R^e,
\]

choose $x_e=0$, then solve recursively:

\[
x_{e-1}=y_{e-1},
\]

\[
x_i=y_i-\theta x_{i+1}
\quad
(e-2\ge i\ge0).
\]

Its kernel is free of rank $1$, generated by

\[
(1,-\theta,\theta^2,\ldots,(-\theta)^e).
\]

Therefore

\[
|\ker L_e(\theta)|=|R|.
\]

The evaluated skew-symmetric block has matrix

\[
\begin{pmatrix}
0&L_e(\theta)^{\mathsf T}\\
-L_e(\theta)&0
\end{pmatrix}.
\]

Its radical is

\[
\ker L_e(\theta)\oplus\ker L_e(\theta)^{\mathsf T}.
\]

Since $L_e(\theta)$ is split surjective, its dual map is injective. Hence

\[
\ker L_e(\theta)^{\mathsf T}=0.
\]

The radical therefore has cardinality

\[
|R|.
\]

Thus its $k$-dimension is

\[
\dim_k R
=
\deg(\eta)f.
\]

\end{proof}

\begin{corollary}[Total Artin radical formula]

Let the pencil have singular multiplicities $m_e$, and let

\[
M_{\mathrm{sing}}=\sum_e m_e
\]

be the total number of singular Kronecker blocks. Let its regular divisor be

\[
\mathcal D_{\mathrm{reg}}
=
\sum_{\xi,e}m_{\xi,e}[\xi,e].
\]

For the closed-point Artin probe ($\eta,f$), the total radical dimension is

\[
\boxed{
\operatorname{rad}_{\eta,f}
M_{\mathrm{sing}}\deg(\eta)f
+
2\deg(\eta)\sum_e m_{\eta,e}\min(e,f).
}
\]

\end{corollary}

\begin{proof}

The radical dimension is additive over the canonical direct-sum decomposition. Singular blocks contribute

\[
\deg(\eta)f
\]

each by Lemma 7.1. Regular blocks contribute

\[
2\deg(\eta)\min(e,f)
\]

when their closed point is $\eta$, and $0$ otherwise, by Theorem 6.1.

\end{proof}

Since the singular multiplicities $m_e$ are recovered from the chain probes, $M_{\mathrm{sing}}$ is known. Therefore define

\[
S_{\eta,f}
=
\frac{
\operatorname{rad}_{\eta,f}
M_{\mathrm{sing}}\deg(\eta)f
}{
2\deg(\eta)
}.
\]

Then

\[
\boxed{
S_{\eta,f}
\sum_e m_{\eta,e}\min(e,f).
}
\]

Taking finite differences gives

\[
S_{\eta,f}-S_{\eta,f-1}
=
\sum_{e\ge f}m_{\eta,e},
\]

and therefore

\[
\boxed{
m_{\eta,f}
2S_{\eta,f}-S_{\eta,f-1}-S_{\eta,f+1}.
}
\]

Thus the regular elementary-divisor multiplicities at every closed point $\eta$ are recovered.

\section{Word realization of Artin probes}

We explain why the Artin radical dimensions above are observable from commutator-word distributions.

Let

\[
R=R_{\eta,f}=K_\eta[\pi]/(\pi^f).
\]

Choose a $k$-basis

\[
\rho_1,\ldots,\rho_d
\]

of $R$, where

\[
d=f\deg(\eta).
\]

Write multiplication in $R$ as

\[
\rho_i\rho_j=\sum_\ell c_{ij}^{\ell}\rho_\ell.
\]

An $R$-valued variable

\[
u\in V\otimes_k R
\]

is represented by $d$ ordinary $V$-variables

\[
u_1,\ldots,u_d
\]

through

\[
u=\sum_i u_i\otimes\rho_i.
\]

The expression

\[
B_R(u,v)-\theta A_R(u,v)
\]

has $d$ $k$-coordinates, each of which is a finite $k$-linear combination of commutators, with coefficients determined by the structure constants of $R$ and the coordinates of $\theta$. Hence the distribution of this expression is an explicit commutator-word distribution.

Let

\[
\chi:R\to\mathbb C^\times
\]

be a generating additive character. Since $R$ is Frobenius, additive characters are parametrized by $R$-linear functionals. The Fourier coefficient associated to $\chi$ gives

\[
\frac{|\operatorname{rad}(B_R-\theta A_R)|}{|V\otimes R|}.
\]

Thus the radical size, and hence its $k$-dimension, is observable.

This applies to the total pencil. By Sections 6 and 7, the observed Artin radical dimensions recover the regular elementary-divisor multiplicities after subtracting the already known singular contribution.

\section{Finiteness of the probe family}

Fix

\[
N=\dim V.
\]

For singular probes, it suffices to use

\[
1\le f\le N+1.
\]

Indeed, a singular block $\mathcal K_e$ has dimension $2e+1$, hence $e\le (N-1)/2$, and the formula

\[
m_e=\delta_{e+1}-2\delta_e+\delta_{e-1}
\]

uses $f\le e+1$.

For regular probes, a block $(\eta,e)$ can occur only if

\[
2e\deg(\eta)\le N.
\]

Therefore one only needs closed points $\eta$ with

\[
\deg(\eta)\le N/2
\]

and exponents

\[
1\le e\le \left\lfloor\frac{N}{2\deg(\eta)}\right\rfloor.
\]

To recover $m_{\eta,e}$ by second differences, use probes up to one additional order:

\[
1\le f\le
\left\lfloor\frac{N}{2\deg(\eta)}\right\rfloor+1.
\]

Since $k$ is finite, there are only finitely many closed points of each bounded degree. Thus, for fixed $k$ and $N$, the total family of probes is finite.

\section{Orbit-coding and the intrinsic projective invariant}

The previous reconstruction is made after choosing a basis of the target $W=k^2$. A different basis acts by

\[
\operatorname{PGL}_2(k)
\]

on the parameter line $\mathbb P^1$. Thus the intrinsic regular datum is not the marked divisor

\[
\mathcal D_{\mathrm{reg}},
\]

but its orbit

\[
\operatorname{PGL}_2(k)\cdot\mathcal D_{\mathrm{reg}}.
\]

We now formalize the finite orbit-coding step.

Let

\[
\mathcal R_{k,N}
\]

be the finite set of decorated regular divisors

\[
\mathcal D=\sum_{\xi,e}m_{\xi,e}[\xi,e]
\]

satisfying

\[
2\sum_{\xi,e}e\deg(\xi)m_{\xi,e}\le N.
\]

Choose an injective code

\[
c:\mathcal R_{k,N}\to\mathbb Z_{\ge0}.
\]

Let

\[
B_0>|\operatorname{PGL}_2(k)|.
\]

Define

\[
\mathcal O_{B_0}(\mathcal D)
=
\sum_{g\in \operatorname{PGL}_2(k)}
B_0^{c(g\mathcal D)}.
\]

Because $B_0>|\operatorname{PGL}_2(k)|$, the base-$B_0$ expansion of $\mathcal O_{B_0}(\mathcal D)$ has no carry ambiguity and determines the multiset

\[
\{c(g\mathcal D):g\in \operatorname{PGL}_2(k)\}.
\]

Since $c$ is injective, this determines the orbit

\[
\operatorname{PGL}_2(k)\cdot\mathcal D.
\]

The singular multiplicities $m_e$ are invariant under target basis change. Thus the pair

\[
\left((m_e)_e,\ \operatorname{PGL}_2(k)\cdot\mathcal D_{\mathrm{reg}}\right)
\]

is the intrinsic projective canonical datum recovered from the word distributions.

\section{Complete pencil theorem}

We now prove the main theorem.

\begin{theorem}[Complete commutator-word tomography of alternating pencils]

Fix

\[
k=\mathbb F_p,\qquad p\neq2,
\]

and

\[
N\ge0.
\]

There exists an explicit finite family of Fourier-tagged commutator-word distributions such that, for every alternating pencil

\[
\beta:\Lambda^2_k V\to k^2,
\qquad
\dim V=N,
\]

the values of these distributions determine the full skew-symmetric Kronecker canonical data of $\beta$, namely:

\begin{itemize}
\item the multiplicities $m_e$ of all singular Kronecker blocks;
\end{itemize}

\begin{itemize}
\item the $\operatorname{PGL}_2(k)$-orbit of the regular decorated divisor
\end{itemize}

\[
\mathcal D_{\mathrm{reg}}\sum_{\xi,e}m_{\xi,e}[\xi,e].
\]

Equivalently, these distributions determine the projective congruence class of the alternating pencil.

\end{theorem}

\begin{proof}

The rectangular Fourier-tagged distributions of Theorem 3.2 with

\[
(C,D)=(F_f,G_f)
\]

recover the ranks

\[
\operatorname{rank}(F_f\otimes A+G_f\otimes B)
\]

for all required $f$. Hence they recover

\[
\delta_f
=
f\dim V-\operatorname{rank}(F_f\otimes A+G_f\otimes B).
\]

By Theorem 5.3,

\[
\delta_f=\sum_e m_e\max(0,f-e),
\]

and finite differences recover every singular multiplicity $m_e$.

Now the singular part is known. For every closed point $\eta$ and every required $f$, the Artin probe of Section 8 recovers the total radical dimension

\[
\operatorname{rad}
=
{\eta,f}.
\]

By Corollary 7.2,

\[
\operatorname{rad}
=
{\eta,f}
M_{\mathrm{sing}}\deg(\eta)f
+
2\deg(\eta)\sum_e m_{\eta,e}\min(e,f).
\]

Since $M_{\mathrm{sing}}$ is known, we compute

\[
S_{\eta,f}
=
\frac{
\operatorname{rad}_{\eta,f}
M_{\mathrm{sing}}\deg(\eta)f
}{
2\deg(\eta)
}
\sum_e m_{\eta,e}\min(e,f).
\]

Finite differences recover every $m_{\eta,e}$. Varying $\eta$ over all closed points of degree at most $N/2$ recovers the marked regular divisor.

Finally, Section 10 converts the marked regular divisor into its intrinsic $\operatorname{PGL}_2(k)$-orbit. Together with the singular multiplicities, this is exactly the skew-symmetric Kronecker canonical datum of the pencil up to projective congruence.

\end{proof}

\section{Consequence for class-two exponent-$p$ groups}

\begin{corollary}[Complete word-moment invariant for groups with $\dim G'=2$]

Let $G$ and $H$ be special class-two exponent-$p$ groups with

\[
\dim_{\mathbb F_p}G'=\dim_{\mathbb F_p}H'=2.
\]

Let

\[
\beta_G:\Lambda^2(G/Z(G))\to G',
\]

\[
\beta_H:\Lambda^2(H/Z(H))\to H'
\]

be their commutator tensors. Then $G\cong H$ if and only if the finite family of Fourier-tagged commutator-word distributions of Theorem 11.1 agrees for $G$ and $H$, after the natural projective target orbit-coding.

\end{corollary}

\begin{proof}

For special class-two exponent-$p$ groups, isomorphism is equivalent to pseudo-isometry of the commutator tensor. When $\dim G'=2$, this tensor is an alternating pencil. By Theorem 11.1, the specified distributions recover the projective congruence class of that pencil, which is precisely its pseudo-isometry class.

\end{proof}

\section{Remarks}

\subsection{Exact distributions, not sampling}

The theorem uses exact finite distributions and Fourier coefficients. It does not address how many random samples are needed to estimate these probabilities.

\subsection{Efficiency}

The orbit-code is finite but may be large. The result is a structural observability theorem, not an optimized group-isomorphism algorithm.

\subsection{Relation to rank-zeta tomography}

Scalar-extension rank tomography recovers rank distributions and zeta functions of rank loci. For pencils, rank-zeta data sees the reduced arithmetic support of the degeneracy divisor but does not by itself recover singular minimal indices or elementary-divisor thicknesses. The present theorem adds chain probes and Artin probes to recover the full canonical data.

\subsection{Relation to word-map fibre enumeration}

The Fourier/rank mechanism is close in spirit to known fibre computations for word maps on class-two exponent-$p$ groups. The new feature here is the design of rectangular and Artin probes that extract the complete Kronecker canonical data of an alternating pencil.


\begin{thebibliography}{99}
\bibitem{ref1} Matthew Levy, Enumerating fibres of commutator words over $p$-groups, arXiv:1602.04093.
\bibitem{ref2} Vladimir V. Sergeichuk, Canonical matrices of forms and pairs of forms over finite and $p$-adic fields, arXiv:1011.3142.
\bibitem{ref3} V. A. Bovdi, T. G. Gerasimova, M. A. Salim, and V. V. Sergeichuk, Reduction of a pair of skew-symmetric matrices to its canonical form under congruence, arXiv:1712.08729.
\bibitem{ref4} James B. Wilson, Decomposing $p$-groups via Jordan algebras, Journal of Algebra 322 (2009), 2642--2679.
\bibitem{ref5} Peter A. Brooksbank and James B. Wilson, Computing isometry groups of Hermitian maps, Transactions of the American Mathematical Society 364 (2012), 1975--1996.
\bibitem{ref6} Peter A. Brooksbank, Yinan Li, Youming Qiao, and James B. Wilson, Improved algorithms for alternating matrix space isometry, ESA 2020.
\bibitem{ref7} Xiaorui Sun, Faster Isomorphism for $p$-Groups of Class 2 and Exponent $p$, arXiv:2303.15412.
\bibitem{ref8} F. R. Gantmacher, The Theory of Matrices, Chelsea, 1959.
\bibitem{ref9} P. Lancaster and M. Tismenetsky, The Theory of Matrices, Academic Press, 1985.
\bibitem{ref10} Thomas Honold, Characterization of finite Frobenius rings, Archiv der Mathematik 76 (2001), 406--415.
\end{thebibliography}
\end{document}
